\chapter{Conclusiones y Análisis de Resultados}

Se presenta el análisis de los resultados experimentales obtenidos a partir 
de las pruebas de rendimiento del generador de \textit{cartoons}. Se contrastan las diferentes 
implementaciones y se analizan las causas arquitectónicas, de software y de entorno físico que justifican el comportamiento observado en los gráficos.

\section{En el entorno local}

El comportamiento de memoria compartida (OpenMP) en el entorno local es el que muestra una gran mejora en los tiempos con respecto a los demas algoritmos. Se asume esto debido a como no se duplican matrices ni se hacen envíos de mensajes, los threads acceden directamente a la imagen cargada en la RAM. De esta forma se obtiene la gran ventaja de OpenMP sobre MPI en todas las configuraciones y resoluciones.
\begin{itemize}
    \item El techo de rendimiento se encuentra en este caso en los 4 threads, y cuando se hace el filtro $5\times 5$, creemos que esto es así por que se tiene el triple de multiplicaciones por pixel. Entonces el speedup mejora porque la carga computacional justifica dividir el trabajo en threads
\end{itemize}

En los algoritmos de MPI y el hibrido consideramos que si bien los resultados mejoran los tiempos del secuencial, estas mejoras son menores que en OpenMP por el tiempo de overhead que se tiene por el pasaje de mensajes. Esto es tiempo en el que OpenMP gana considerable ventaja. 
    

\section{En el entorno de el cluster}

El comportamiento de la implementación de memoria compartida (OpenMP) en el cluster muestra que el speedup permanece practicamente plano (oscilando entre $0.91\times$ y $1.15\times$) sin importar que la cantidad de threads aumente de 2 a 32. 
Se asume que tal comportamiento se atribuye a dos factores:
En los entornos de cluster (como SLURM), cuando se ejecuta un proceso multiprocesador de forma nativa sin configurar directivas de asignación específicas (por ejemplo, omitiendo la reserva de múltiples núcleos virtuales por tarea), el sistema operativo limita por afinidad (\textit{affinity binding}) el proceso completo a \textit{un único núcleo físico de CPU}. 
Como consecuencia, aunque la directiva OpenMP en software cree 4, 8, 16 o 32 threads, todos estos threads son forzados por el kernel del sistema operativo a alternarse en el mismo procesador. Esto no solo impide la computación paralela real, sino que añade sobrecarga debido al intercambio de contexto (*context switching*) y a la invalidación de caché de datos.
    
El diseño del algoritmo divide la carga mediante \texttt{\#pragma omp parallel sections} asignando un thread al posterizado y otro al borroneado. Dado que la convolución bidimensional del borroneado (filtro de blur y Sobel) requiere miles de operaciones aritméticas en punto flotante por píxel, esta sección tarda considerablemente más tiempo que el posterizado (que consiste en una división simple por píxel). 
Al ser la sección de borroneado el cuello de botella masivo, y al ejecutarse esta de forma secuencial debido a la restricción de afinidad sobre un solo procesador,  el speedup general queda limitado por la Ley de Amdahl a un máximo de $\approx 1.15\times$, haciendo que la adición de threads no tenga impacto real.

A diferencia de OpenMP, las implementaciones que utilizan pasaje de mensajes (MPI) y la version hibrida logran mostrar una \textit{escalabilidad positiva}:
Al invocar el programa mediante \texttt{mpirun -n p}, se crean $p$ procesos independientes del sistema operativo; asi, el planificador del sistema y el entorno MPI distribuyen estos procesos en diferentes nucleos fisicos.
En cuanto a speedup se observa que tanto para MPI como para la configuración híbrida, el mismo aumenta comforme la imagen es más grande.
